\section{Ambientes Complejos}
En esta sección se involucra la \textbf{relajación de restricciones}, donde hay presentes \emph{problemas, algoritmos, estrategias y problemas no deterministicos}
\begin{definition}[Problemas de Busqueda]
    \begin{enumerate}
        \item Completamente observables
        \item deterministicos
        \item Estaticos
        \item ambientes conocidos
    \end{enumerate}
\end{definition}


\begin{definition}[Problemas discretos de hallar un buen estado]
    \begin{enumerate}
        \item Busqueda hill-climbing
        \item simulated anealing
        \item local beam search
        \item evolutionary algorithms
    \end{enumerate}
\end{definition}

Hasta ahora los algoritmos de busqueda (BFS, DFS, UCS, A*) asumen ambientes \textbf{completamente observables}, 
\textbf{deterministicos}, \textbf{estaticos} y \textbf{conocidos}, y su objetivo es encontrar un \emph{camino} desde el estado inicial hasta la meta.

Sin embargo, hay clases de problemas donde lo que importa no es el camino sino el \textbf{estado final} en si mismo. esta sección 
estudia tres familias de técnicas para esos casos:
\begin{enumerate}
    \item \textbf{Busqueda local en espacios discretos:} Hill-climbing, simulated annealing, local beam search y algortimos evolutivos.
    \item \textbf{Busqueda local en espacios continuos:} discretización, gradiente empirico, newton-raphson y programacion lineal
    \item \textbf{Busqueda con acciones no deterministas:} arboles AND-OR y planes condicionales
\end{enumerate}

En busqueda local \textbf{no se guarda el camino recorrido}; el estado actual es todo lo que importa. Esto reduce dramaticamente el uso de memoria, al costo
de no garantizar optimalidad global.

\section{Busqueda Local y Problemas de optimización}
\begin{definition}[busqueda local]
    Clase de algoritmos que dado un estado actual, generan sus vecinos y se mueven al mejor de ellso segun una \textbf{funcion objetivo}
    $f: \mathcal{S} \to \mathbb{R}$. no mantienen frontera ni arbol de busqueda completo.
\end{definition}

El espacio de busqueda se puede visualizar como un \emph{paisaje de estados:} los estados son posiciones horizontales y su valor $f$ es la altura. El objetivo tipico
es encontrar el \textbf{maximo global} o minimo dependiendo de la formulacion.

\subsection{Hill-Climbing}
\begin{definition}
    Mantiene registro de un unico estado actual, y en cada iteración, se desplaza al vecino con \textbf{mayor valor} de $f$. termina cuando ningun vecino supera al estado actual. 
    Tambien llamado \emph{steepest-ascent} p \emph{busqueda local voraz}
\end{definition}

\begin{example}[Hill-Climbing clasico]
    \begin{enumerate}
        \item Estado actual $\leftarrow$ estado inicial
        \item Generar todos los vecinos del estado actual
        \item Si algun vecino tiene $f$ mayor: moverse al mejor vecino y volver a 2
        \item Si ningun vecino mejora: \textbf{terminar} (maximo local)
    \end{enumerate}
\end{example}

